\documentclass[a4paper,12pt]{article}
\usepackage{color}
\usepackage{graphicx, gensymb}
\usepackage{amsfonts, cancel, amssymb}
\usepackage{amsmath, amsthm, array, esvect, esint}
\usepackage{typearea}
\usepackage[a4paper,left=2cm,right=2cm,top=2cm,bottom=3cm,bindingoffset=5mm]{geometry}
\newcommand\numberthis{\addtocounter{equation}{1}\tag{\theequation}}
\newcommand{\bra}{}
\newcommand{\kom}{}
\newcommand{\brak}{}
\newcommand{\braket}{}
\newcommand{\intx}{}
\newcommand{\intr}{}
\newcommand{\kla}{}
\newcommand{\klam}{}
\newcommand{\klammer}{}
\newcommand{\oper}{}
\newcommand{\tecxt}{}
\newcommand{\Bra}{}
\newcommand{\Ket}{}

\begin{document}

\title{Derivation of the exact electromagnetic field of a moving point charge}
\author{Joachim Schwardt}
\date{\today}
\maketitle

\section{Decoupling the \textsc{Maxwell}-equations}
The \textsc{Maxwell}-equations in a vacuum with natural units ($\hbar=c=1$) are:
\begin{align*}
\nabla\cdot E&=\mu_0\varrho&\nabla\times E&=-\partial_tB\\
\nabla\cdot B&=0&\nabla\times B&=\mu_0j+\partial_tE
\end{align*}
Relativistic equivalent:
\begin{align*}
F^{\mu\nu}&=\begin{pmatrix}
0&-E \\ E & -\epsilon_{\mu\nu k}B_k
\end{pmatrix}&\widetilde{F}^{\mu\nu}&=\begin{pmatrix}
0&-B \\ B & \epsilon_{\mu\nu k}E_k
\end{pmatrix} \\
\partial_\mu F^{\mu\nu}&=\mu_0J^\nu&\partial_\mu\widetilde{F}^{\mu\nu}&=0
\end{align*}
Define the vector potentials and the \textsc{Lorentz}-gauge as follows:
\begin{align*}
E&=:-\nabla\phi-\partial_t A&B&=:\nabla\times A &\nabla A&=-\partial_t\phi
\end{align*}
Relativistic equivalent is $F^{\mu\nu}=\partial^\mu A^\nu -\partial^\nu A^\mu$ and $\partial_\mu A^\mu =0$.\\
Using these potentials and the gauge for the \textsc{Maxwell}-equations gives:
\begin{align*}
-\nabla^2\phi -\partial_t\nabla A&=\mu_0\varrho&\nabla(\nabla A)-\nabla^2 A+\nabla\partial_t\phi+\partial_t^2A&=\mu_0j\\
\Rightarrow\ \ \kla\left( {\nabla^2-\partial_t^2} \right)\phi&=-\mu_0\varrho &\Rightarrow\ \ \kla\left( {\nabla^2-\partial_t^2} \right)A&=-\mu_0j
\end{align*}
Relativistic equivalent is $\partial_\mu\partial^\mu A^\nu =\mu_0J^\nu$.\\
We now want to solve these equations using \textsc{Green}-functions. Hence, we need to find a solution to the PDE\ \ $\square G(r,t)=\delta(r)\delta(t)$. Taking the \textsc{Fourier}-transform gives $\kla\left( {\omega^2-k^2} \right)G(k,\omega)=\frac{1}{4\pi^2}$. The following convention was applied:
\begin{align*}
\mathcal{F}\klam\left[ {f} \right]&=\frac{1}{(2\pi)^{n/2}}\intr\int_{\mathbb{R}^{n}}{f(r)e^{-i\brak\langle {k}| {r}\rangle}}\,\tecxt\text{d}^nr&\mathcal{F}^{-1}\klam\left[ {f} \right]&=\frac{1}{(2\pi)^{n/2}}\intr\int_{\mathbb{R}^{n}}{f(k)e^{i\brak\langle {k}| {r}\rangle}}\,\mathrm{d}^{n}k
\end{align*} 
This equation can be solved by the inverse transformation. The residual theorem is used to solve the integral:
\begin{align*}
G(r,t)&=\frac{1}{16\pi^4}\intr\int_{\mathbb{R}^{3}}{\intr\int_{\mathbb{R}^{ }}{\frac{e^{i\kla\left( {k\cdot r-\omega t} \right)}}{\omega^2-k^2}}\,\mathrm{d}^{ }\omega}\,\mathrm{d}^{3}k= \\
&=-\frac{1}{16\pi^4}\intr\int_{\mathbb{R}^{3}}e^{ik\cdot r}\lim_{\epsilon\rightarrow 0}{\intx\int_{\pi}^{0}\frac{e^{itk+it\epsilon  e^{iz}}}{\epsilon e^{iz}\kla\left( {-2k+\epsilon e^{iz}} \right)}\cdot\epsilon ie^{iz}+\frac{e^{-itk-it\epsilon e^{iz}}}{\epsilon e^{iz}\kla\left( {2k+\epsilon e^{iz}} \right)}\cdot\epsilon ie^{iz}\,\mathrm{d}z}\,\mathrm{d}^{3}k=\\
&=\frac{i}{32\pi^4}\intr\int_{\mathbb{R}^{3}}{\frac{e^{ik\cdot r}}{k}\intx\int_{\pi}^{0}e^{itk}-e^{-itk}\,\mathrm{d}z}\,\mathrm{d}^{3}k=\\
&=\frac{1}{32i\pi^3}\intx\int_{0}^{\infty}e^{itk}-e^{-itk}\intx\int_{0}^{2\pi}\intx\int_{0}^{\pi}\frac{e^{ikr\cos\vartheta}}{k}\cdot k^2\sin\vartheta\,\mathrm{d}\vartheta\mathrm{d}\varphi\mathrm{d}k=\\
&=\frac{1}{16i\pi^2}\intx\int_{0}^{\infty}k\kla\left( {e^{itk}-e^{-itk}} \right)\frac{e^{ikr\cos\vartheta}}{-ikr}\bigg\vert_{0}^\pi\,\mathrm{d}k=\\
&=\frac{1}{16\pi^2r}\intx\int_{0}^{\infty}e^{itk-ikr}-e^{itk+ikr}-e^{-itk-ikr}+e^{-itk+ikr}\,\mathrm{d}k=\\
&=\frac{1}{8\pi^2r}\intx\int_{0}^{\infty}\cos\kla\left( {k(t-r)} \right)-\cos\kla\left( {k(t+r)} \right)\,\mathrm{d}k= \\
&=\frac{1}{8\pi^2r}\lim_{R\rightarrow\infty}\klam\left[ {\frac{\sin\kla\left( {R(t-r)} \right)}{t-r}-\frac{\sin\kla\left( {R(t+r)} \right)}{t+r}} \right]=:\frac{1}{8\pi^2r}\lim_{R\rightarrow\infty}I_R^{-r}-I_R^{+r}\\
I_R^{-r}\klam\left[ {f} \right]&=\intr\int_{\mathbb{R}^{ }}{\frac{\sin\kla\left( {Rt} \right)}{t}f(t+r)}\,\mathrm{d}^{ }t=\intr\int_{\mathbb{R}^{ }}{\frac{\sin x}{x}}f\kla\left( {r+\frac{x}{R}} \right)\,\mathrm{d}^{ }x\overset{R\rightarrow\infty}{\longrightarrow}\pi f(r)\\
G(r,t)&=\frac{1}{8\pi r}\kla\left( {\delta(t-r)-\delta(t+r)} \right)\overset{?}{=}-\frac{\delta (t\pm r)}{4\pi r}
\end{align*}
\section{Solving the \textsc{Maxwell}-equations}
A general solution to the equations can be found using the obtained \textsc{Green}-function:
\begin{align*}
\phi&=\frac{\mu_0}{4\pi}\intr\int_{\mathbb{R}^{3}}{\intr\int_{\mathbb{R}^{ }}{\frac{\varrho(r',t')\delta(t-t'-|r-r'|)}{|r-r'|}}\,\mathrm{d}^{ }t'}\,\mathrm{d}^{3}r'\\
A&=\frac{\mu_0}{4\pi}\intr\int_{\mathbb{R}^{3}}{\intr\int_{\mathbb{R}^{ }}{\frac{j(r',t')\delta(t-t'-|r-r'|)}{|r-r'|}}\,\mathrm{d}^{ }t'}\,\mathrm{d}^{3}r'
\end{align*}
Let's calculated these integrals for the motion of a point charge along some arbitrary curve $R(t)$. We get $\varrho=q\delta (r-R(t))$ and $j=q\dot{R}(t)\delta(r-R(t))$. Therefor:
\begin{align*}
\phi&=\frac{q\mu_0}{4\pi}\intr\int_{\mathbb{R}^{ }}{\intr\int_{\mathbb{R}^{3}}{\frac{\delta(r'-R(t'))\delta(t-t'-|r-r'|)}{|r-r'|}}\,\mathrm{d}^{3}r'}\,\mathrm{d}^{ }t'=\\
&=\frac{q\mu_0}{4\pi}\intr\int_{\mathbb{R}^{ }}{\frac{\delta(t-t'-|r-R(t')|)}{|r-R(t')|}}\,\mathrm{d}^{ }t'=:\frac{q\mu_0}{4\pi}\intr\int_{\mathbb{R}^{ }}{\frac{\delta(f(t'))}{|r-R(t')|}}\,\mathrm{d}^{ }t'
\end{align*}
Since $f'(t)=-1+\frac{r-R(t)}{|r-R(t)|}\cdot\dot{R}(t)=:-1+n\cdot\beta<0$ and $f(t\rightarrow\infty)<0$, $f(t\rightarrow -\infty)>0$ we have exactly one root at $t'=:t_r$ which is obtained from solving the implicit equation $t_r=t-|r-R(t_r)|=:t-d$. Hence, $\delta(f(t'))=\frac{\delta(t'-t_r)}{|f'(t_r)|}=\frac{\delta(t'-t_r)}{1-\beta\cdot n}$:
\begin{align*}
\phi&=\frac{q\mu_0}{4\pi d}\frac{1}{1-\beta\cdot n}&A&=\beta\phi
\end{align*}
\section{Deriving the electromagnetic field}


\end{document}